\subsection{MATRIX OF A SEMI-LINEAR MAPPING}

Let A, B be two rings, \(\sigma: \mathbf{A} \to \mathbf{B}\) a homomorphism of A into B, E a right (resp. left) A-module with basis \((e_i)_{i \in I}\) and F a right (resp. left) B-module with basis \((f_k)_{k\in K}\) . Let \(u:\mathbf{E}\to \mathbf{F}\) be a semi-linear mapping relative to \(\sigma\) and \(u(e_i) = \sum_{k\in K}f_ku_{ki}\left(\text{resp.} u(e_i) = \sum_{k\in K}u_{ki}f_k\right)\) , where the \(u_{ki}\) are therefore elements of B; by definition, the matrix \(M(u) = (u_{ki})\) of type \(\mathbf{K}\times \mathbf{I}\) is also called the matrix of \(u\) with respect to the bases \((e_i)\) and \((f_k)\) . By the same calculation as in Proposition 2 of no. 4, it is immediately verified that for all \(x\in \mathbf{E}\) , if I and K are finite,

\[
M (u (x)) = M (u). \sigma (M (x))\tag{26 D}
\]

(resp.

\[
{ } ^ { t } M ( u ( x ) ) = \sigma ( { } ^ { t } M ( x ) ) . { } ^ { t } M ( u ) ) .\tag{26 G}
\]

Let C be a third ring, \(\tau: \mathbf{B} \to \mathbf{C}\) a homomorphism, G a right (resp. left) C-module with basis \((g_l)_{l \in \mathbf{L}}\) and \(v\) a semi-linear mapping of F into G relative to \(\tau\) ; if \(M(v)\) is the matrix of \(v\) with respect to \((f_k)\) and \((g_l)\) and \(M(v \circ u)\) the matrix of \(v \circ u\) relative to \((e_i)\) and \((g_l)\) , then, if I, K and L are finite,

\[
M (v \circ u) = M (v). \tau (M (u))\tag{27 D}
\]

(resp.

\[
{ } ^ { t } M ( v \circ u ) = \tau ( { } ^ { t } M ( u ) ) . { } ^ { t } M ( v ) ) .\tag{27 G}
\]

To show for example (27 D), note that for all \(x \in E\) , by (26 D),

\[
\begin{array}{r l} M (v \circ u) \cdot \tau (\sigma (M (x))) & = M (v (u (x))) \\ & = M (v) \cdot \tau (M (u (x))) = M (v) \cdot \tau (M (u)) \cdot \tau (\sigma (M (x))), \end{array}
\]

whence (27 D) by Proposition 1 of no. 3.

Suppose finally that \(\sigma: \mathbf{A} \to \mathbf{B}\) is an isomorphism; then recall that \(^t u: \mathbf{F}^* \to \mathbf{E}^*\) is a semi-linear mapping relative to \(\sigma^{-1}\) (§ 2, no. 5); when I and K are finite, the matrix \(^t u\) with respect to the dual bases \((f_k^*)\) and \((e_i^*)\) is given by

\[
M (^ {t} u) = \sigma^ {- 1} (^ {t} M (u))\tag{28}
\]

for, by definition, supposing for example that E and F are right modules, \(\langle^t u(f_k^*), e_i \rangle^\sigma = \langle f_k^*, u(e_i) \rangle\) when \(\sigma\) is denoted by \(x \mapsto x^\sigma\) .

Remark. Let A be a ring and \(\sigma\) an anti-endomorphism of A (no. 3); consider the two following situations:

(GD) E is a left A-module, F a right A-module and u a Z-linear mapping of E into F such that \(u(ax) = u(x)\sigma(a)\) for \(a \in A\) , \(x \in E\) ; in other words, u is a semi-linear mapping relative to \(\sigma\) of the right \(A^{0}\)-module E into the right A-module F.

(DG) E is a right A-module, F a left A-module and u a Z-linear mapping of E into F such that \(u(xa) = \sigma(a)u(x)\) for \(a \in A\) , \(x \in E\) ; in other words, u is a semi-linear mapping relative to \(\sigma\) of the left \(A^{0}\)-module E into the left A-module F.

In the two cases, the matrix \(M(u)\) of u relative to bases of E and F has its elements in A; if these bases are finite, then, for all \(x \in E\) , we have the respective formulae

\[
M (u (x)) = M (u). \sigma (M (x))\tag{17 GD}
\]

\[
{ } ^ { t } M ( u ( x ) ) = \sigma ( { } ^ { t } M ( x ) ) . { } ^ { t } M ( u ) ,\tag{17 DG}
\]

the products on the two sides being calculated in A. This follows immediately from (26 D) and (26 G) respectively.

